\clearpage
\section{Effective Addressing Modes}\label{page:effective-addressing-modes}

An effective-address field is an operand descriptor with an instruction-defined role and interpretation width. A
value role reads or writes the selected register, immediate, or memory value. An address role uses a register or
immediate as an address value and uses a memory-form EA as the calculated address without an initial data read. A
control-target role uses a register or immediate value directly and reads an interpretation-width target from a
memory form. Memory forms continue into segment pre-translation and paging.

The seven-bit compact EA field is embedded in the opcode payload. Any displacement, absolute address, immediate, or
EXT0 descriptor bytes are appended as EA payload bytes. When an instruction contains multiple payload-bearing
operands, payloads are consumed in instruction operand order.

Multi-byte scalar EA payloads, such as displacements, absolute addresses, and immediates, are little-endian. Payload
names ending in s are signed and are sign-extended before use; payload names without s are unsigned unless the
consuming instruction explicitly says otherwise. The decoder consumes EXT0 descriptor bytes in the displayed stream
order, byte 0 followed by byte 1, and bit numbering restarts at 7 for each byte.

Indexed scale is the EA interpretation width declared by the consuming form. B, W, L, and Q widths select scales 1,
2, 4, and 8. A form whose width is \texttt{operation\_size} uses its selected instruction size. LEA and SEGLEA use
their suffix; size-less address and control-target forms use Q.

For every compact or EXT0 PC-based EA, the PC term is the architectural PC naming byte 0 of the current instruction.
The PC term remains that instruction-start value throughout operand evaluation.

\clearpage
\subsection{Compact EA Addressing Modes}
Compact EA forms encode direct Rn and SP operands, common Rn/SP/PC memory
operands, absolute memory operands, immediate operands, and the EXT0 escape.

Compact memory forms that do not name SP or PC use the operation default data segment. SP memory forms are fixed to
SS, and PC memory forms are fixed to CS.

\begin{manualformblock}{3.65in}
\textbf{\texttt{Register Direct}}\par
\vspace{2pt}
\begin{tabularx}{\linewidth}{@{}p{1.10in}X@{}}
\textbf{Assembler syntax} & Rn(r), SP\\
\textbf{EA encoding} & 000rrrr selects Rn(r); 1101000 selects SP.\\
\textbf{EA type} & register operand, not a memory reference.\\
\textbf{Generation} & The operand value is the selected register value. No effective memory address is generated.\\
\textbf{Segment} & No segment is selected.\\
\textbf{Payload} & No appended EA payload bytes.\\
\textbf{Update} & No auto-update.\\
\end{tabularx}\par
\manualeadirectflow{Direct register operand}{REGISTER}{Rn or SP contents}{REGISTER VALUE}
\end{manualformblock}
\begin{manualformblock}{5.07in}
\textbf{\texttt{Rn Memory}}\par
\vspace{2pt}
\begin{tabularx}{\linewidth}{@{}p{1.10in}X@{}}
\textbf{Assembler syntax} & [Rn(r)], [Rn(r) + disp8s], [Rn(r) + disp16s], [Rn(r) + disp32s], [Rn(r) + disp64]\\
\textbf{EA encoding} & 001rrrr has no displacement; 010rrrr, 011rrrr, 100rrrr, and 101rrrr select disp8s, disp16s, disp32s, and disp64.\\
\textbf{EA type} & memory operand.\\
\textbf{Generation} & offset = temporary Rn(r) + displacement. The no-displacement form uses displacement 0.\\
\textbf{Segment} & Uses the operation default data segment unless the instruction defines a different default.\\
\textbf{Payload} & Only displacement forms append payload bytes. Displacement payload is little-endian and has the size named by the EA form.\\
\textbf{Update} & No auto-update in compact Rn memory forms.\\
\end{tabularx}\par
\manualeaadditivememoryflow{Rn memory address generation}{BASE REGISTER}{Rn(r)}{optional displacement}{EFFECTIVE ADDRESS}
\end{manualformblock}
\begin{manualformblock}{5.07in}
\textbf{\texttt{SP Memory}}\par
\vspace{2pt}
\begin{tabularx}{\linewidth}{@{}p{1.10in}X@{}}
\textbf{Assembler syntax} & [SP], [SP + disp8s], [SP + disp16s], [SP + disp32s], [SP + disp64]\\
\textbf{EA encoding} & 1101001 has no displacement; 1100000, 1100001, 1100010, and 1100011 select disp8s, disp16s, disp32s, and disp64.\\
\textbf{EA type} & memory operand.\\
\textbf{Generation} & offset = temporary SP + displacement. The no-displacement form uses displacement 0.\\
\textbf{Segment} & Fixed to SS; no segment field is encoded.\\
\textbf{Payload} & Only displacement forms append payload bytes. Displacement payload is little-endian and has the size named by the EA form.\\
\textbf{Update} & No auto-update in compact SP memory forms.\\
\end{tabularx}\par
\manualeaadditivememoryflow{SP memory address generation}{STACK POINTER}{SP}{optional displacement}{EFFECTIVE ADDRESS}
\end{manualformblock}
\begin{manualformblock}{5.07in}
\textbf{\texttt{PC Memory}}\par
\vspace{2pt}
\begin{tabularx}{\linewidth}{@{}p{1.10in}X@{}}
\textbf{Assembler syntax} & [PC + disp8s], [PC + disp16s], [PC + disp32s], [PC + disp64]\\
\textbf{EA encoding} & 1100100, 1100101, 1100110, and 1100111 select disp8s, disp16s, disp32s, and disp64.\\
\textbf{EA type} & memory operand.\\
\textbf{Generation} & offset = temporary PC term + displacement. The PC term is the PC value supplied to operand evaluation by the execution model.\\
\textbf{Segment} & Fixed to CS; no segment field is encoded.\\
\textbf{Payload} & A displacement payload is always present. Payload is little-endian and has the size named by the EA form.\\
\textbf{Update} & No auto-update for PC forms.\\
\end{tabularx}\par
\manualeaadditivememoryflow{PC-relative address generation}{PROGRAM COUNTER}{PC term}{displacement}{EFFECTIVE ADDRESS}
\end{manualformblock}
\begin{manualformblock}{4.25in}
\textbf{\texttt{Absolute Memory}}\par
\vspace{2pt}
\begin{tabularx}{\linewidth}{@{}p{1.10in}X@{}}
\textbf{Assembler syntax} & [abs32s], [abs64]\\
\textbf{EA encoding} & 1101010 selects abs32s; 1101011 selects abs64.\\
\textbf{EA type} & memory operand.\\
\textbf{Generation} & offset = sign-extended abs32s or unsigned abs64.\\
\textbf{Segment} & Uses the operation default data segment unless the instruction defines a different default.\\
\textbf{Payload} & The absolute address payload is little-endian and has the size named by the EA form.\\
\textbf{Update} & No auto-update.\\
\end{tabularx}\par
\manualeasimplememoryflow{Absolute memory address generation}{ABSOLUTE ADDRESS}{address payload}{EFFECTIVE ADDRESS}
\end{manualformblock}
\begin{manualformblock}{3.97in}
\textbf{\texttt{Immediate}}\par
\vspace{2pt}
\begin{tabularx}{\linewidth}{@{}p{1.10in}X@{}}
\textbf{Assembler syntax} & imm8s, imm16s, imm32s, imm64\\
\textbf{EA encoding} & 1101100, 1101101, 1101110, and 1101111 select imm8s, imm16s, imm32s, and imm64.\\
\textbf{EA type} & immediate operand, not a memory reference.\\
\textbf{Generation} & The operand value is decoded from the immediate payload. Signed payloads are sign-extended according to the consuming instruction's operand rule.\\
\textbf{Segment} & No segment is selected.\\
\textbf{Payload} & The immediate payload is little-endian and has the size named by the EA form.\\
\textbf{Update} & No auto-update. Immediate forms are not valid destinations unless an instruction explicitly defines such a form.\\
\end{tabularx}\par
\manualeaimmediateflow{Immediate operand generation}{payload value}{IMMEDIATE VALUE}
\end{manualformblock}
\begin{manualformblock}{2.35in}
\textbf{\texttt{EXT0 Escape}}\par
\vspace{2pt}
\begin{tabularx}{\linewidth}{@{}p{1.10in}X@{}}
\textbf{Assembler syntax} & EXT0, EXT0/disp8s, EXT0/disp16s, EXT0/disp32s, EXT0/disp64\\
\textbf{EA encoding} & 1110100 selects no displacement; 1110000, 1110001, 1110010, and 1110011 select disp8s, disp16s, disp32s, and disp64.\\
\textbf{EA type} & memory operand described by the appended EXT0 descriptor.\\
\textbf{Generation} & The EXT0 descriptor selects segment, base, index, and auto-update behavior. The compact EA escape selects only the displacement size.\\
\textbf{Segment} & EXT0 has no implicit segment. Segment-qualified forms encode SEG(s); SP and PC indexed forms use SS and CS; default base-update forms use the operation default data segment.\\
\textbf{Payload} & Payload order is EXT0 descriptor byte or bytes first, then the optional displacement payload selected by the compact EA escape.\\
\textbf{Update} & Only EXT0 forms with ++ or {-}{-} update the temporary operand-evaluation image.\\
\end{tabularx}\par
\end{manualformblock}


@COMPACT_EA_TABLE@
\label{table:compact-ea-encoding}

\clearpage
\subsection{Extended EA Descriptor}
A compact EXT0 EA value selects an extended descriptor and the displacement width. The descriptor bytes then select
segment, base, index, zero-base, SP/PC indexed, and auto-update behavior.

Each EXT0 descriptor selects an SREG SEG(s), fixed SS or CS through SP or PC, or the operation default data segment.

\begin{manuallistedbitdiagram}{Compact Effective-Address Field}
\manualbitfieldrow{EA[6:0]}{%
\manualbitlabel{6}
\manualbitlabel{4}
\manualbitlabel{3}
\manualbitlabel{0}
}{%
\manualbitfieldtext{mode}{3}
\manualbitfieldtext{register/form}{4}
}
\end{manuallistedbitdiagram}

\begin{manuallistedbitdiagram}{EXT0 One-Byte Descriptor Forms}
\manualbitrow{SEG + base}{%
\manualbitfixed{0}{1}
\manualbitfieldcode{s}{3}
\manualbitfieldcode{b}{4}
}
\manualbitrow{SEG + zero}{%
\manualbitfixed{1}{1}
\manualbitfieldcode{s}{3}
\manualbitfixed{0011}{4}
}
\manualbitrow{default base++}{%
\manualbitfixed{1}{1}
\manualbitfieldcode{b}{4}
\manualbitfixed{100}{3}
}
\manualbitrow{default {-}{-}base}{%
\manualbitfixed{1}{1}
\manualbitfieldcode{b}{4}
\manualbitfixed{101}{3}
}
\end{manuallistedbitdiagram}

\begin{manuallistedbitdiagram}{EXT0 Two-Byte Indexed Descriptor Byte Layouts}
\manualbitfieldrow{SEG + base + index}{%
\manualbytepairlabels
}{%
\manualbitfixed{1}{1}
\manualbitfieldcode{s}{3}
\manualbitfieldcode{mode}{4}
\manualbitgap{1}
\manualbitfieldcode{b}{4}
\manualbitfieldcode{i}{4}
}
\manualbitfieldrow{SP/PC + index}{%
\manualbytepairlabels
}{%
\manualbitfixed{1000101}{7}
\manualbitfieldcode{x}{1}
\manualbitgap{1}
\manualbitfieldcode{mode}{4}
\manualbitfieldcode{i}{4}
}
\end{manuallistedbitdiagram}


\manualtablecaption{EA Payload Names}
\begin{manuallongtable}{@{}>{\raggedright\arraybackslash}p{0.95in}>{\raggedright\arraybackslash}p{0.45in}>{\raggedright\arraybackslash}p{0.70in}>{\raggedright\arraybackslash}p{3.30in}@{}}
\toprule
\rowcolor{ManualHeaderFill}
\textbf{Payload} & \textbf{Bytes} & \textbf{Value} & \textbf{Use}\\
\midrule
\endfirsthead
\multicolumn{4}{l}{\scriptsize\itshape Table \themanualtable\ (continued)}\\
\toprule
\rowcolor{ManualHeaderFill}
\textbf{Payload} & \textbf{Bytes} & \textbf{Value} & \textbf{Use}\\
\midrule
\endhead
@EA_PAYLOAD_ROWS@
\bottomrule
\end{manuallongtable}

For a compact displacement, absolute address, or immediate, the selected payload follows the EA field at that
operand's payload position. EXT0 begins with the compact escape value and one or two descriptor bytes; a displaced
EXT0 form places its selected displacement after those descriptor bytes. When an instruction has more than one
payload-bearing EA operand, the decoder consumes each complete EA payload sequence in instruction operand order.

@EXT0_SECTION@

@AUTO_UPDATE_SECTION@
